% OWNER: sections/s2_1_transport.tex. Edit only this file.
\seclabel
\claim{\secnum{2.1}{sec:transport}\quad Reusing one HTTP/2 connection cuts request time from 106.5\ms{} to 16.0\ms}

\lead{The site serves HTTP/2 over TLS 1.3 and does not advertise HTTP/3.} The TLS handshake selects HTTP/2 through ALPN, with the cipher suite \code{TLS_AES_256_GCM_SHA384} and a chain of 4 certificates from Let's Encrypt. Responses carry no \code{alt-svc} header, the header through which a server advertises HTTP/3 (QUIC). We infer the absence of HTTP/3 from that missing header alone; we did not probe UDP port 443.

\lead{An nginx front end passes requests to an ASP.NET application.} nginx's own HTML error page appears when a request body is malformed (Section~\ref{sec:auth}); the application's error messages and its SignalR hub identify the framework. The ASP.NET defaults cited in Sections~\ref{sec:auth} and~\ref{sec:feed} rest on this identification.

\lead{REST requests share one HTTP/2 connection, and the live feed needs a second connection (\figref{layers}).} A REST request travels as one HTTP/2 stream carrying the \code{Authorization: Bearer} header; later requests open new streams on the same connection and skip DNS, TCP and TLS. The live feed starts with an HTTP/2 POST to negotiate, then opens a second connection and upgrades it with an HTTP/1.1 \code{Upgrade: websocket} request. The server forces this second connection: its HTTP/2 SETTINGS frame lacks \code{ENABLE_CONNECT_PROTOCOL}, the setting a server must send before a client may carry a WebSocket over HTTP/2 [1]. Every client, a browser included, therefore opens a separate connection for the feed.

% OWNER: figs/layers.tex. Edit only this file.
\fig[wide]{layers}{REST requests share one HTTP/2 connection; the live feed needs a second, WebSocket connection}[Protocol stack per channel, client side, from the application down to the connection set-up]{%
\begin{tikzpicture}[x=1mm,y=1mm,
  box/.style={draw=navy,line width=0.5pt,fill=white,minimum width=30mm,minimum height=8mm,inner sep=0pt,font=\fontsize{7.5}{9}\selectfont,anchor=north west},
  lname/.style={anchor=west,align=left,font=\fontsize{7.5}{9}\selectfont\bfseries\color{navy}},
  note/.style={anchor=west,align=left,text width=69mm,inner sep=0pt,font=\fontsize{7}{8.4}\selectfont\color{greyD}},
  hd/.style={anchor=south,align=center,font=\fontsize{7.5}{9}\selectfont\bfseries\color{navy}}]
% column heads
\node[hd] at (58,-1.2) {REST\\[-0.5mm]{\fontsize{7}{8.4}\selectfont\mdseries\color{greyD}\code{/api}}};
\node[hd] at (92,-1.2) {Live feed\\[-0.5mm]{\fontsize{7}{8.4}\selectfont\mdseries\color{greyD}\code{/hubs/main}}};
% stacks: rows of 8 mm from y = -2
\foreach \y/\n/\a/\b in {-2/Application/{REST API}/SignalR, -10/Framing/{HTTP/2 stream}/WebSocket, -18/HTTP/{HTTP/2}/{HTTP/1.1}, -26/TLS/{TLS 1.3}/{TLS 1.3}, -34/TCP/TCP/TCP}{
  \node[box] at (43,\y) {\a};
  \node[box] at (77,\y) {\b};
  \node[lname] at (23,\y-4) {\n};
}
% DNS, before connecting
\node[box,draw=greyD,dashed,minimum width=64mm] at (43,-45) {DNS};
\node[lname,text=greyD] at (23,-49) {Before\\[-0.6mm]connecting};
% bracket over TLS, TCP and DNS of both stacks
\draw[navy,line width=0.6pt] (21.5,-26.3) -- (20,-26.3) -- (20,-52.7) -- (21.5,-52.7);
\draw[navy,line width=0.6pt] (20,-39.5) -- (18.5,-39.5);
\node[anchor=east,align=right,font=\fontsize{7}{8.4}\selectfont\bfseries\color{navy}] at (18,-39.5) {paid again\\on every new\\connection};
% short notes, one per row
\node[note] at (111,-6)  {REST: Bearer token on every request.\\Feed: token checked once, when the connection opens};
\node[note] at (111,-14) {REST: one stream per request, many per connection.\\Feed: JSON records, each ended by 0x1E};
\node[note] at (111,-22) {REST: chosen by ALPN. Feed: HTTP/2 POST to negotiate,\\then \code{Upgrade: websocket} on its own connection};
\node[note] at (111,-30) {1 round trip; the feed runs its own handshake};
\node[note] at (111,-38) {REST: closed by the server after 15 to 30\,s idle.\\Feed: kept open by a request every 3\,s};
\node[note] at (111,-49) {\code{contest.timefolio.net} $\rightarrow$ 221.168.38.40};
\end{tikzpicture}
}{Each box is one protocol; both channels start from the same DNS lookup. Later REST requests open new HTTP/2 streams on the same connection and skip DNS, TCP and TLS. The live feed needs the second connection because the server's HTTP/2 SETTINGS frame lacks \code{ENABLE_CONNECT_PROTOCOL}, so WebSockets over HTTP/2 [1] are unavailable to any client, a browser included. On the server, an nginx front end passes requests to an ASP.NET application. The SETTINGS frame was read with Python's \code{h2} library; the front end is identified by its error page, the application by its error messages.}{DART Quant Team}


\lead{A new connection costs 106.5\ms{} and a reused one 16.0\ms, so set-up takes about 90\ms{} of the first request.} We sent 5 requests to the same endpoint (\code{Portfolio/Violations}) back to back in one curl run, about 170\ms{} per run, and repeated the run 10 times at 1\,s intervals. The first request in each run pays for the DNS lookup, the TCP connection and the TLS handshake; the other 4 reuse the connection (\figref{timing}). The stage medians add up to 106.0\ms; the median of the totals, 106.5\ms, comes from different runs.

% OWNER: figs/timing.tex. Edit only this file.

\fig[wide]{timing}{A new connection costs 106.5\ms, a reused one 16.0\ms}[One curl request to \code{Portfolio/Violations}, medians of 10 runs (40 reused requests); time runs down, same scale in both panels]{%
\begin{tikzpicture}[x=1mm,y=1mm,
  f/.style={font=\fontsize{7}{8.2}\selectfont},
  host/.style={font=\fontsize{7}{8.2}\selectfont\bfseries,text=navy,anchor=south,align=center},
  tl/.style={greyM,line width=0.6pt,dotted,-{Stealth[length=1.4mm]}},
  msg/.style={-{Stealth[length=1.6mm]},line width=0.7pt,navy},
  rep/.style={-{Stealth[length=1.6mm]},line width=0.7pt,navyM},
  ml/.style={f,text=navy,sloped,above,inner sep=0.6pt},
  stg/.style={line width=0.5pt,navy},
  sl/.style={f,anchor=east,align=right,text=navy},
  ptitle/.style={font=\fontsize{7.5}{9}\selectfont\bfseries,text=navy,anchor=west}]
\def\tT#1{-0.72*(#1)}% time in ms -> y in mm (0.72 mm per ms, both panels)
% ============================================ panel (a): new connection
\def\Rx{0}\def\Cx{22}\def\Sx{58}
\node[ptitle] at (-24,14) {(a) New connection: DNS, TCP and TLS come before the request};
\node[host] at (\Rx,2.5) {DNS\\resolver};
\node[host] at (\Cx,2.5) {Client\\(curl)};
\node[host] at (\Sx,2.5) {Server\\(nginx)};
\foreach \x in {\Rx,\Cx,\Sx} {\draw[tl] (\x,1) -- (\x,{\tT{110}});}
% DNS: the split inside the 49.3 ms is not measured; arrows drawn one-way 3 ms
\draw[msg] (\Cx,{\tT{0}}) -- node[ml]{A query} (\Rx,{\tT{3}});
\fill[navyXL] (\Rx-1,{\tT{3}}) rectangle (\Rx+1,{\tT{46.3}});
\draw[rep] (\Rx,{\tT{46.3}}) -- node[ml]{A reply} (\Cx,{\tT{49.3}});
% TCP: one RTT
\draw[msg] (\Cx,{\tT{49.3}}) -- node[ml]{SYN} (\Sx,{\tT{54.85}});
\draw[rep] (\Sx,{\tT{54.85}}) -- node[ml,pos=0.78,fill=white,inner sep=0.6pt]{SYN-ACK} (\Cx,{\tT{60.4}});
% TLS 1.3: one RTT, then the rest of curl's TLS time on the client
\draw[msg] (\Cx,{\tT{60.4}}) -- node[ml,pos=0.66]{ACK, ClientHello} (\Sx,{\tT{65.95}});
\draw[rep] (\Sx,{\tT{65.95}}) -- node[ml,below,fill=white,inner sep=0.6pt,pos=0.35]{ServerHello ... Finished} (\Cx,{\tT{71.5}});
\fill[navyXL] (\Cx-1,{\tT{71.5}}) rectangle (\Cx+1,{\tT{89.3}});
\node[f,anchor=west,text=greyD,align=left] at (\Cx+1.8,{\tT{81.4}}) {rest of curl's TLS time};
% request to first byte
\draw[msg] (\Cx,{\tT{89.3}}) -- node[ml]{Finished, \code{HEADERS} (request)} (\Sx,{\tT{94.85}});
\fill[navyXL] (\Sx-1,{\tT{94.85}}) rectangle (\Sx+1,{\tT{100.45}});
\node[f,anchor=west,text=greyD] at (\Sx+1.8,{\tT{97.65}}) {server};
\draw[rep] (\Sx,{\tT{100.45}}) -- node[ml,below,inner sep=0.4pt]{\code{HEADERS} 200, \code{DATA}} (\Cx,{\tT{106}});
% stage brackets (left)
\draw[stg] (-3+1,{\tT{0}-0.2}) -- ++(-1,0) -- (-3,{\tT{49.3}+0.2}) -- ++(1,0);
\node[sl] at (-4.6,{\tT{24.65}}) {DNS lookup\\\textbf{49.3}\ms\\{\color{greyD}no DNS cache}\\{\color{greyD}in WSL}};
\draw[stg] (-3+1,{\tT{49.3}-0.2}) -- ++(-1,0) -- (-3,{\tT{60.4}+0.2}) -- ++(1,0);
\node[sl] at (-4.6,{\tT{54.85}}) {TCP \textbf{11.1}\ms\\= 1 RTT};
\draw[stg] (-3+1,{\tT{60.4}-0.2}) -- ++(-1,0) -- (-3,{\tT{89.3}+0.2}) -- ++(1,0);
\node[sl] at (-4.6,{\tT{74.85}}) {TLS 1.3 \textbf{28.9}\ms\\{\color{greyD}handshake 1 RTT,}\\{\color{greyD}12.4\ms{} in Python}};
\draw[stg] (-3+1,{\tT{89.3}-0.2}) -- ++(-1,0) -- (-3,{\tT{106}+0.2}) -- ++(1,0);
\node[sl] at (-4.6,{\tT{97.65}}) {request to first\\byte \textbf{16.7}\ms};
\draw[stg] (\Sx+12,{\tT{0}}) -- ++(1,0) -- (\Sx+13,{\tT{106}}) -- ++(-1,0);
\node[f,text=navy,rotate=-90,anchor=south] at (\Sx+14.5,{\tT{53}}) {total \textbf{106.5}\ms{} (median of totals; stage medians sum to 106.0\ms)};
% ============================================ panel (b): reused connection
\def\Cb{94}\def\Sb{126}
\node[ptitle] at (\Cb-8,14) {(b) Reused connection: the request only};
\node[host] at (\Cb,2.5) {Client\\(curl)};
\node[host] at (\Sb,2.5) {Server\\(nginx)};
\foreach \x in {\Cb,\Sb} {\draw[tl] (\x,1) -- (\x,{\tT{24}});}
\draw[msg] (\Cb,{\tT{0}}) -- node[ml]{\code{HEADERS} (request)} (\Sb,{\tT{5.55}});
\fill[navyXL] (\Sb-1,{\tT{5.55}}) rectangle (\Sb+1,{\tT{10.45}});
\draw[rep] (\Sb,{\tT{10.45}}) -- node[ml,below,inner sep=0.4pt]{\code{HEADERS} 200, \code{DATA}} (\Cb,{\tT{16}});
\draw[stg] (\Cb-3+1,{\tT{0}-0.2}) -- ++(-1,0) -- (\Cb-3,{\tT{16}+0.2}) -- ++(1,0);
\node[sl] at (\Cb-4.6,{\tT{8}}) {\textbf{16.0}\ms};
\node[f,anchor=north west,align=left,text=navy] at (\Cb-6,{\tT{24}-2}) {1 RTT (11.1\ms) plus about 5\ms{} on the server:\\a new HTTP/2 stream on the open connection};
% run-to-run spread strip: min to max over 10 runs (40 for reused), tick = median; 0.18 mm per ms
\begin{scope}[shift={(\Cb+14,{\tT{24}-17})}]
\def\s{0.16}
\node[f,anchor=south west,font=\fontsize{7}{8.2}\selectfont\bfseries,text=navy] at (-20,1.6) {Spread over runs, min to max (ms)};
\foreach \lab/\lo/\md/\hi/\r in {DNS/36.6/49.3/94.1/0, TCP/9.4/11.1/12.3/1, TLS/17.8/28.9/40.4/2, To first byte/15.7/16.7/19.8/3, {Total, new}/80.5/106.5/156.9/4, {Total, reused}/14.8/16.0/20.3/5} {
  \pgfmathsetmacro{\yy}{-3.9*\r}
  \node[f,anchor=east,text=navy] at (-0.8,\yy) {\lab};
  \draw[navyL,line width=1.8pt] ({\s*\lo},\yy) -- ({\s*\hi},\yy);
  \draw[navy,line width=1pt] ({\s*\md},\yy-1.3) -- ({\s*\md},\yy+1.3);
  \node[f,anchor=west,text=greyD] at ({\s*160+0.8},\yy) {\lo{} to \hi};
}
\draw[greyM,line width=0.4pt] (0,-21.5) -- ({\s*160},-21.5);
\foreach \t in {0,50,100,150} {\draw[greyM,line width=0.4pt] ({\s*\t},-21.5) -- ({\s*\t},-22.3); \node[f,anchor=north,text=greyD] at ({\s*\t},-22.3) {\t};}
\end{scope}
% legend and scale bar
\begin{scope}[shift={(-24,{\tT{110}-6})}]
\draw[msg] (0,0) -- (7,0); \node[f,anchor=west] at (8,0) {client to server};
\draw[rep] (32,0) -- (39,0); \node[f,anchor=west] at (40,0) {reply};
\fill[navyXL] (54,-1.2) rectangle (56,1.2); \node[f,anchor=west] at (57,0) {local work};
\draw[greyD,line width=0.6pt] (78,{-0.72*5}) -- (78,{0.72*5});
\foreach \t in {-5,5} {\draw[greyD,line width=0.6pt] (77.2,{0.72*\t}) -- (78.8,{0.72*\t});}
\node[f,anchor=west,text=greyD] at (79.5,0) {10\ms{} on both panels};
\end{scope}
\end{tikzpicture}%
}{One-way times are drawn as half the 11.1\ms{} TCP round trip (RTT). Only the stage boundaries are measured: curl's \code{time_*} values are cumulative, and a stage is the difference between adjacent values; ``request to first byte'' runs from the end of curl's TLS stage to the first byte of the reply. The split inside the DNS lookup and inside curl's TLS stage is not measured. Labels give medians; the strip gives the range over the 10 runs (40 reused requests), with the median as a tick. All 50 requests returned HTTP 200. 06:03 KST, 24 Nov 2026. Time-space drawing after Kurose and Ross [6].}{DART Quant Team}

% OWNER: figs/pytimings.tex. Edit only this file.
\fig[wide]{pytimings}{The Python client pays 1.2 to 1.3\,s once per process and 31 to 39\ms{} per reused request}[Time of one REST request by client, moment and connection state, 24 Nov 2026]{%
{\tabfont\begin{tabularx}{\linewidth}{>{\RaggedRight\arraybackslash}p{40mm} >{\RaggedRight\arraybackslash}p{42mm} R{17mm} @{\hspace{6pt}}>{\RaggedRight\arraybackslash}X}
\thd{Client, moment} & \thd{Request} & \thd{Time} & \thd{Where the time goes}\\
curl, 06:03 (before the open) & Reused connection & 16.0\ms & 1 RTT (11.1\ms) plus about 5\ms{} on the server (\figref{timing})\\\rowsep
\hl Python, during the session & Reused connection & 31 to 39\ms & Cause of the gap to curl not separated: client, time of day, market-hours load\\\rowsep
Python, fresh process & First request & 1.2 to 1.3\,s & Module loading and SSL set-up, once per process; the virtual environment lies on the Windows drive\\\rowsep
Python, same process & First request of a second client & 130\ms & New connection only; the one-time cost is already paid\\
\end{tabularx}}}{The shaded row is the cost of every wrapper request after start-up. The 1.2 to 1.3\,s is paid once by \code{warm_up()} at start, before the first order. Python times come from the \code{transport} live checks (\figref{tests}).}{DART Quant Team}


\lead{\figref{timing} uses only a repeat run with a valid token; an earlier run with an expired token gave reused requests as fast as 12.2\ms{} and was excluded.} That run had not recorded status codes, and 12.2\ms{} fits a server that rejects the token before touching any data. The repeat run recorded a status code on every request.


\lead{A reused request costs one network round trip of 11.1\ms{} plus about 5\ms{} at the two ends.} curl's TCP stage runs from sending the SYN segment to receiving the SYN-ACK, which is one round trip between client and server. Its median was 11.1\ms, and its standard deviation across the 10 runs was 0.9\ms, so the network path was stable. A request on a reused connection also needs one round trip, request out and reply back. Of the 16.0\ms, 11.1\ms{} is therefore network time and about 5\ms{} is spent in the server and the client together; we did not separate the two. Faster client code can save only part of those 5\ms.

\lead{Two stages of the new connection reflect our client set-up more than the server.} Under WSL, every new connection paid 37 to 94\ms{} for its DNS lookup in the curl runs (median 49.3\ms), and a separate Python probe measured 30 to 39\ms{} per lookup. Lookups 1\,s apart did not get faster, so no cache on our path answered them; a system with a caching resolver would pay the lookup once. The wrapper still connects by host name, since the certificate names \code{contest.timefolio.net}: connecting by IP address would save the lookup and lose hostname verification. curl's 28.9\ms{} for TLS also includes curl's own TLS start-up. Timed in Python on an open TCP connection, the TLS 1.3 handshake itself took 12.4\ms, one round trip.

\lead{From the Python client, the first request of a process took 1.2 to 1.3\,s and a reused request 31 to 39\ms{} (\figref{pytimings}).} The reused figure, measured during the trading session, is about twice curl's 16.0\ms{} measured before the open. The two measurements differ in client, time of day and market-hours load, and we have not separated these causes. Almost all of the 1.2 to 1.3\,s is one-time module loading and SSL set-up (the virtual environment lies on the Windows drive): a second client created in the same process paid 130\ms{} for its new connection. The one-time cost therefore comes with the process, and a new connection from Python costs about 130\ms. At start-up the wrapper calls \code{warm_up()}, which sends one cheap GET to \code{Portfolio/Session} (the order window, Section~\ref{sec:queries}) and ignores its 401, so the one-time cost is paid before the first order (D18).

\lead{One trial per interval puts the server's idle cut-off between 15 and 30\,s.} httpx drops a connection after 5\,s idle by default [4], which hid the server's behaviour in our first try. With that limit raised to 600\,s, a connection was reused after 5\,s and 15\,s of silence and replaced after 30\,s. Each interval was tried once, so the cut-off may differ between connections. An HTTP/2 server normally announces an idle close with a GOAWAY frame, and the client then opens a new connection for its next request; we did not capture frames, so this mechanism is assumed rather than observed. A request meets a closed socket only when it is sent while the GOAWAY is still in flight. A GET caught in that race is sent once more, but an order the server may have received is never resent (Section~\ref{sec:orders}), so the wrapper closes idle connections itself before the server does.

\verdict{the wrapper keeps one \code{httpx.AsyncClient(http2=True)} for the life of the process. It sets \code{keepalive_expiry} to 10\,s, below 15\,s, the longest idle time after which a connection was still reused, so the wrapper closes an idle connection before the server can (D1, D2). A request within 10\,s of the previous one reuses the connection; after a longer gap the wrapper opens a new one, about 130\ms{} from Python. Before an order, the \code{ValidateProduct} GET of the pre-order check (Section~\ref{sec:guidelines}) reopens the connection, so \code{AddOrder} itself goes out on an open connection.}
